\scanpage{25}
\mspage{21}
% archive sheet number: 24

et
\[
\unsure{M|\mathcal M}(J,\omega_1)=\{\omega_1,J_{<\omega_1}\}
\]
magasin qui consiste en \uline{deux} éléments seulement, et \uline{c'est une moquette}. Ici, on a
\[
|J_{<\omega_1}|^{\circ}=\varnothing,\quad\text{mais bien sûr}\quad J^{\circ}_{<\omega_1}\neq\varnothing.
\]

\deleted{Choix dual} Traitement dual pour
\[
M(I,\partial I)\setminus\{\omega_1\}\overset{\mathrm{def}}=M(J,\omega_0)
\]
où $\omega_0$ est donné en \uline{tant que} plus petit élément de $J\equiv I\setminus\{\omega_1\unsure{)|\}}$. Cette fois, c'est les éléments
\[
J_{>a}\qquad(a\in J)
\]
qu'il faut introduire, en plus des ``lieux'' $a\in J$ et des ``intervalles finis'' \unsure{de $J$}
\[
J_{a,b}.
\]

Prenons enfin le cas du magasin résiduel de $\partial I$ dans $M$
\[
M(I,\partial I)\setminus\partial I\overset{\mathrm{def}}=M(J)\qquad J=I\setminus\partial I
\]
Cette fois, il faut introduire à la fois les
\[
J_{<a}\quad\text{et les}\quad J_{>a}\qquad(a\in J),
\]
% slanted
\margin{\unsure{ceci}, y compris si $a$ est minimal ou maximal dans $J$!}
pour les ``traces'' des intervalles \deleted{\unsure{en}} $I_{\omega_0,a},I_{a,\omega_1}$, et \deleted{de plus} enfin, de plus, un objet noté
\[
[\unsure{\mathbb J|J}]
\]
qui tient lieu de $I_{\omega_0,\omega_1}$. Les \unsure{introductions}
